% Methods 2.1-2.9 prose — PARKED 2026-11-01, NOT the current draft.
%
% ⚠⚠ DYLAN PULLED THIS OUT OF THE SCAFFOLD on 2026-11-01 and reverted Methods to its
% fourteen one-line items (doc rev 917). Two reasons, in his words: "you're adding
% way too much detail that I don't know if we need ... generally trying to include
% way too many of our findings from the project that were just for our own
% understanding", and "we are introducing things in the wrong order for all of this".
% 2.1's second paragraph was named as the specific example of unreadable.
%
% ⚠ IT IS KEPT FOR SENTENCES, NOT AS A DRAFT. Nothing here goes back into the paper
% until the section ORDER is settled in the scaffold, and then only the parts a
% reader needs in order to follow the method. The verified figures below are still
% good; the prose around them is what was wrong.
%
% ⚠⚠ ONE PARAGRAPH IN 2.1 IS ALSO FACTUALLY WRONG, found 2026-11-01 while answering
% Dylan's comment on it. It says the chain excludes positions where no donor
% haplotype has a called allele. That is true of exactly ONE position on chr21,
% where the chain excludes 8,328. The real exclusion is a single 8,327-position
% block (10,326,676 to 12,968,320, ~2.6 Mb) of zero genetic distance that no single
% donor haplotype spans, so the copying model has no path through it — decision K's
% untraversable block, recorded in data/cohort/chr21/in_chain.npy's own provenance.
% Do not reuse that paragraph; the fact belongs in one sentence of the dataset
% description, which Methods does not yet have at all.
%
% ⚠ THE DATASET DESCRIPTION IS THE REBUILD'S FIRST ITEM. Dylan: "We need a
% description of the actual dataset we are using." 2.1 gave a position axis without
% ever saying whose genomes these are, which is part of why the order read wrong.
%
% Every figure below was verified against source or artifacts this session, not
% quoted from a docstring:
%   chr21 site axis 340,849 / chain T = 332,521  <- data/cohort/chr21/in_chain.npy
%   genome-wide T = 22,134,562 over 22 autosomes <- sum over in_chain.npy
%   K = 88 from 44 individuals                   <- haplotypes.npy shape (88, 340849)
%   weights = "uniform"                          <- scripts/arm2_attack.py:204,
%                                                   scripts/arm1_attack.py:94
%   diploid tilt = eta/2 per haplotype           <- src/mechanism/ffbs.py:617
%   N_e = 10,000, r = 1.26                       <- src/mechanism/hmm.py NE, R
%   chr21 eta_tau 75.35 (0.55) to 517.17 (0.80)  <- arm2_chr21_kl_verified.json,
%     so epsilon_tau = 2*eta_tau is 150.7..1034.3   refined block, 10 targets
%   genome eta_tau 3308.5..3716.8 at 0.80        <- logs/B_acc_arm2/genome_sweep/*.json
%     epsilon_tau 6617..7434, log10(1-tau)          aggregate, acc_ceiling 0.80
%     -2873 (HG00096) to -3228 (NA19350)
%   Pinsker/Gaussian excess ratio 1.2554 /       <- recomputed from KL = 2(2a-1)^2;
%     1.2721 / 1.3289 at 0.55 / 0.65 / 0.80        sqrt(pi/2) = 1.2533 is the limit
%   envelope: D1 > every real decoy in 180/180  <- logs/B_acc_arm2/arm2_chr21_ceiling
%     chr21 arm 2 cells, worst ratio 0.478403,     .json (MEASURED THIS SESSION; the
%     worst real ceiling 0.7075 at the 0.80        instrument check reproduced the
%     label; 264/264 arm 1 chr21 and 60/60         verified D1 kl on all 60 cells to
%     genome for the conditions measured there     6.82e-13)
%
% OPEN DECISIONS, all Dylan's:
%   A. tau IS EXACTLY 1.0 IN FLOAT64 PAST eta ~ 19, and the reported grid runs
%      eta ~ 75 to 517 on chr21. So at every reported operating point the
%      theorem's bound TV <= tau is VACUOUS, not merely loose. The Introduction
%      as approved states the guarantee in tau. 2.9 and 4.2 own the ladder, but a
%      referee will reach this before Section 2.9. DECIDED 2026-10-31 (Dylan):
%      "leave 1.4 alone". The Introduction carries no caveat, so 2.3 below is
%      where a reader first meets the fact, and the sentence in 2.3 that says it
%      is load-bearing rather than incidental. The counter-argument -- that a
%      referee gets there before 2.9 -- was logged and overruled.
%   B. The Introduction says a new genome "can be mapped or projected onto the
%      graph"; 2.1 says we JOIN a phased, already-called sample column. Those are
%      consistent (the Introduction states the general capability), but if the
%      general claim is not wanted, 1.3's first sentence is where to narrow it.
%   C. u_ceiling and the "30.749104% of beta_v the target cannot call" figure are
%      WEIGHTING-SPECIFIC and predate the 2026-10-24 uniform ruling. Not quoted
%      here. They need re-measuring under uniform before any section uses them.
%   D. src/mechanism/hmm.py documents four deviations between PanMixer's shipped
%      sampler and its own published Methods, with a measured consequence. That
%      is 4.4's material if it is used at all, not 2.4's. Not stated here.
%   E. THE CALIBRATION IS ALREADY IN 2.3. eta_tau = arctanh(tau)/delta_u is stated
%      in 2.3 above, where the mechanism is defined. The scaffold's 2.11 line
%      assigns it to 2.11 instead, "because it is a protocol step and not a
%      relationship between bounds". Both placements are defensible and only one
%      should carry it. 2.6 to 2.9 below deliberately do NOT restate it, so the
%      choice stays open: either 2.11 points back to 2.3, or the formula moves
%      from 2.3 to 2.11 and 2.3 names only delta_u. Dylan's call.
%   F. 2.9's envelope paragraph reports a result measured this session, not an
%      argument. It replaces the one open measurement the suite had left: the
%      extremal decoy's dominance over D4_farthest was UNRESOLVED by Monte Carlo
%      at the 0.55, 0.60 and 0.65 labels (the sampled estimator's median per-cell
%      standard error there was 1.007x the exact divergence itself). The exact
%      computation settles all 60 cells with zero sampling error.
%   G. THE RANKING of the three real decoy conditions is NOT in 2.9 and should not
%      be: on exact divergence the medians order D4 > D3 > D2 at all six labels,
%      and D4 exceeds D2 in exactly the 8 of 10 targets that have a real
%      concordance gap and in neither of the 2 that do not. That belongs with the
%      decoy conditions in 2.12, and 2.12 still says "three conditions" and must
%      be brought to four.

\section{Methods}

\subsection{Fixed graph and mapped paths}

The graph is fixed and public, and nothing in this paper modifies it: the haplotype
paths already embedded in it are the donor panel, and a release is a new path over
the same coordinates. Positions are the variant records of the graph's own call set,
and the subset the mechanism runs over---the chain---is every position at which at
least one donor haplotype has a called allele. That exclusion is a function of the
panel alone and is therefore independent of any target: where no donor has a called
allele the sampler has nothing to emit. Each chain position additionally carries a
genetic position in centiMorgans interpolated from a genetic map, and that every
chain position carries a genuine interpolated value is a checked invariant rather
than an assumption. On chromosome~21 the site axis holds $340{,}849$ records, of
which $T = 332{,}521$ are in the chain; genome-wide, across the 22 autosomes,
$T = 22{,}134{,}562$. The donor panel is $K = 88$ haplotype paths from 44
individuals.

A target genome enters as a path over the same axis. Our target inputs are phased
and called against the same reference build as the graph, so a sample column already
is a path once joined to the site axis; the join is strict on position, reference
allele and the full alternate-allele string, which is what makes the allele indices
transferable---under a relaxed join, allele~2 in the target and allele~2 in the panel
need not be the same piece of sequence. No read-alignment step is involved. A
position with no record in the target's call set is marked uncalled and contributes
nothing to the utility, so the attainable utility is strictly below one for every
target.

One privacy parameter covers one genome. A release spanning several chromosomes
splits a single budget across them with per-chromosome shares summing to one, and
the single-chromosome case is the whole budget on one chromosome; spending the full
budget on each of 22 chromosomes independently would be 22 releases of one person
rather than one release, and the shares are checked to sum to one wherever the whole
set is visible. Because the divergence the budget is denominated in is additive
across chromosomes under that split (Section~\ref{sec:exact}), a curator who changes
the scope changes the number of positions, and therefore the per-position budget,
but still receives the ceiling they asked for.

\subsection{Two target regimes}

The mechanism and the guarantee are the same in both regimes; what differs is what
the graph already contains about the target. No quantity in this paper is reported
without naming which regime it belongs to.

In the first, the target is external: its genome contributed to neither the graph nor
the donor panel, and it is absent from every panel an attack draws on. Nothing is
held out, the full panel is available to the baseline, and nothing in the mechanism
is a function of the target except the tilt. This is the regime in which we report
genome-wide results.

In the second, the target is one of the individuals the graph was built from. This is
the regime of the head-to-head against PanMixer, and it requires per-target
leave-one-out: both of the target's haplotype paths are removed from the donor panel,
leaving $K = 86$, so that the support and the baseline remain independent of the
target. Positions that only the target called leave the chain with them, so each
target in this regime has its own chain, its own baseline, its own utility weights
and its own ceiling, and quantities from different targets are consequently not
pooled. Results in this regime are reported for chromosome~21 alone.

Leave-one-out removes the target from the panel but not from the graph: the topology
and the site axis were derived with that individual's assemblies among the inputs.
The guarantee in this regime is therefore conditional on the graph, and we state that
rather than claim it away.

\subsection{Cohort-tilted private path release}

Given the graph, the donor panel and a target path $p$, the mechanism releases a path
$y$ drawn from
\begin{equation}
Q^{\tau}_{p}(y) \;=\; \frac{R_D(y)\,\exp\!\big(\eta_\tau\, u(p,y)\big)}{Z_p},
\end{equation}
where $R_D$ is a target-independent distribution over the paths the donor panel can
produce, $u$ is a bounded measure of agreement with the target, and $\eta_\tau$ sets
how strongly the baseline is tilted toward it. The support is $R_D$'s own support, so
no output the panel cannot produce is ever released and no output it can produce is
excluded; both the support and the baseline are built from the graph and the panel
alone.

The utility is a weighted average of per-position indicators,
\begin{equation}
u(p,y) \;=\; \sum_{v} \beta_v \,\mathbf{1}[\,y_v = p_v\,],
\qquad \sum_{v} \beta_v = 1,
\end{equation}
so $u \in [0,1]$ by construction, and that it does is asserted at run time rather
than argued. The weights are uniform over the chain, $\beta_v = 1/T$.

Uniform weighting is a choice and is argued for rather than inherited. Under it, $u$
has a one-sentence meaning---the fraction of the target's called positions that the
release reproduces---which is comparable across targets, panels and tools. The
alternative we began with, $\beta_v \propto 1/\mathrm{support}(v)$, weights a
property of our data rather than of the position: it directs the budget at the
positions where our own panel knows least (genome-wide, the $75{,}998$ chain
positions with support below 40---$0.343\%$ of the chain---carry $2.491\%$ of the
weight, and a position supported by a single donor carries $85.5$ times the mean
weight), which makes the released object a function of our call rate and not
reproducible by anyone holding a different panel. Uniform weighting is also the
neutral instrument for a claim about the mechanism, since concentrating the weights
is itself a lever on the exchange rate between privacy and utility, and it is the
least concentrated weighting we considered: $\sum_v \beta_v^2$ is
$4.5178\times10^{-8}$ against inverse-support's $7.3440\times10^{-8}$, so reporting
under uniform weights is conservative in that exchange rate by about
$\sqrt{1.63} = 1.27$.

The tilt magnitude is calibrated as $\eta_\tau = \operatorname{arctanh}(\tau)/\Delta
u$, where $\Delta u$ is the range of the utility being tilted with. Here $\Delta u =
1$ by construction, so $\eta_\tau = \operatorname{arctanh}(\tau)$; the range is
carried explicitly rather than assumed, because an unnormalised utility does not
weaken the guarantee but removes it.

The budget does not thin as the path grows. Because the weights sum to one, the total
tilt between a path matching the target everywhere and one matching it nowhere is
$\exp(\eta_\tau)$ however the path is partitioned: each position moves less and there
are proportionally more of them. What does change with $T$ is numerical headroom.

Zero tilt is not a special case. At $\eta_\tau = 0$ every tilt potential is one and
the draw is a sample from $R_D$ alone---a donor mosaic carrying no target information
whatever---reached by the identical code path, which is why that column cannot drift
away from the mechanism it is the floor of.

The sampler takes either $\tau$ or $\eta_\tau$ and never both, the two being related
by $\tau = \tanh(\eta_\tau)$. Past $\eta_\tau \approx 19$ that expression is exactly
$1$ in double precision, so at the tilt magnitudes we report the equivalent $\tau$ is
$1$ and $\eta_\tau$ is the parameter that remains exact.

\subsection{HMM-based implementation}

$R_D$ is a Li--Stephens haplotype-copying chain over the donor panel
\cite{li2003}. The chain visits the chain positions in order; at each one it occupies
a single donor haplotype and emits the allele that donor carries there; and the
genetic distance between consecutive positions decides how readily it switches donor.
The transition between consecutive positions is
\begin{equation}
A_t \;=\; (1 - s_t)\,I \;+\; \frac{s_t}{K}\,J,
\end{equation}
with $s_t$ the \emph{total} switch mass, so the per-state coefficient is $s_t/K$ and
the probability of actually leaving a given donor is $s_t(K-1)/K$. The switch mass
follows the closed form $p_{\mathrm{stay}} = e^{-d/K} + (1 - e^{-d/K})/K$ with
$d = \Delta x\, N_e\, r$, where $\Delta x$ is the genetic distance in centiMorgans,
$N_e = 10{,}000$ and $r = 1.26$; the implementation computes the decomposition and
the closed form independently and checks that they agree.

Missingness enters as an indicator potential and not as a renormalised transition
row. $R_D$ is the copying chain conditioned on never occupying a donor that has no
called allele at the position it occupies, and the conditioning is applied as a
$0/1$ factor that the forward pass normalises away. Masking a transition row and
renormalising it is not well defined here: across a step of zero genetic distance the
transition is the identity, so a donor available at one position and unavailable at
the next has a masked row summing to exactly zero.

A release is drawn by forward filtering followed by backward sampling, which draws
exactly from $Q^{\tau}_{p}$ rather than approximating it. Every forward column is
retained in double precision, because backward sampling needs the whole array and not
the running vector a forward-only simulation keeps. The transition matrix is never
materialised: its rank-one-plus-diagonal structure makes each forward step linear in
$K$ from a single scalar, where a dense matrix per position would be about
$21$\,GB at chromosome-21 scale. The release is sampled and never maximised.
Maximum-a-posteriori decoding returns the single most probable state sequence, which
is not a draw from the release distribution and carries no guarantee at all while
producing entirely plausible output; a test walks the package's syntax tree and fails
on any use of \texttt{argmax}.

The log normaliser is accumulated as the sum of the forward scaling constants less
the baseline's own, since the missingness conditioning contributes to both. It must
lie in $[0, \eta_\tau]$, and checking that it does is the cheapest correctness test
available.

\subsection{Diploid releases}

A diploid target's two haplotype paths are one private input and receive one release,
because that is what the guarantee covers and what one release per genome means. The
two chains are drawn independently, each at half the tilt exponent, and the utility
is the mean over the two haplotypes rather than their sum. The mean is what keeps the
utility inside $[0,1]$, so a single privacy parameter covers the whole diploid
release; the sum would double the range and with it silently double the calibration
the guarantee depends on.

\subsection{The privacy guarantee}

Fix the graph, the donor panel and a privacy parameter $\tau \in [0,1)$. Write
$\mathcal{P}(G)$ for the set of admissible input paths and $Q^{\tau}_{p}$ for the
release distribution the mechanism induces from input $p$. The mechanism is
\emph{$\tau$-private against the cohort} for a target $g$ if
\begin{equation}
\max_{1 \le i \le n} \mathrm{TV}\!\big(Q^{\tau}_{p_g},\, Q^{\tau}_{p_i}\big)
\;\le\; \tau ,
\end{equation}
and \emph{$\tau$-private pairwise} if the stronger
\begin{equation}
\sup_{p,\,q \,\in\, \mathcal{P}(G)}
\mathrm{TV}\!\big(Q^{\tau}_{p},\, Q^{\tau}_{q}\big) \;\le\; \tau
\end{equation}
holds, the supremum running over every admissible pair and not only over cohort
members. The second implies the first for every target at once, and it is the second
that the mechanism satisfies.

Total variation is the right currency because of what it means operationally.
Consider the adversary this paper takes as its reference: it holds one release, both
candidate input paths and the mechanism itself, and must say which of the two
produced the release, under an equal prior over the two. By the Neyman--Pearson
lemma the optimal such test succeeds with probability exactly
\begin{equation}
P^{*}_{\mathrm{succ}}(p,q)
\;=\; \tfrac{1}{2}\big(1 + \mathrm{TV}(Q^{\tau}_{p}, Q^{\tau}_{q})\big),
\label{eq:succ}
\end{equation}
so a pairwise bound of $\tau$ caps optimal identification accuracy at $(1+\tau)/2$
and caps the advantage over guessing at $\tau/2$. The parameter therefore bounds
twice the adversary's advantage, not its success rate. The distinction is worth one
sentence because the other reading is the tempting one and is not merely loose: as a
bound on the success rate it would assert that at $\tau = 0$ no test succeeds, where
a coin succeeds half the time.

For the mechanism of Section~2.3 the guarantee holds pairwise and in two forms. For
any two admissible input paths $p, q \in \mathcal{P}(G)$,
\begin{equation}
\mathrm{TV}\!\big(Q^{\tau}_{p}, Q^{\tau}_{q}\big) \;\le\; \tau,
\qquad\text{and}\qquad
e^{-\varepsilon_\tau} \;\le\; \frac{Q^{\tau}_{p}(y)}{Q^{\tau}_{q}(y)}
\;\le\; e^{\varepsilon_\tau}
\quad\text{for every } y \text{ in the common support},
\label{eq:thm}
\end{equation}
where
\begin{equation}
\varepsilon_\tau \;=\; 2\operatorname{arctanh}(\tau)
\;=\; \log\frac{1+\tau}{1-\tau}.
\end{equation}
The first is an averaged statement about the two distributions. The second is
pointwise, holds output by output, and is the stronger of the two; it is also the
form a reader arriving from differential privacy will look for first.

The pointwise form makes the cost of the reported operating points explicit, and we
state that cost here rather than leave it to be worked out. Because $\Delta u = 1$ we
have $\varepsilon_\tau = 2\eta_\tau$, and the tilt magnitudes we report are large.
Over the ten external targets of the second regime the chromosome-21 grid runs
$\eta_\tau$ from $75.35$ at the $0.55$ operating point to $517.17$ at $0.80$, so
$\varepsilon_\tau$ runs from about $151$ to about $1034$; genome-wide at the $0.80$
point $\varepsilon_\tau$ runs from $6617$ to $7434$ across the same ten targets. At
those magnitudes $\tau = \tanh(\eta_\tau)$ is $1$ to double precision---genome-wide
at $0.80$ it is between $1 - 10^{-2873}$ and $1 - 10^{-3228}$---so at the settings we
report neither the pairwise bound nor the pointwise ratio bound carries any
information. What they carry is scope: they hold for every admissible pair at every
parameter value, which is what makes this a bounded release mechanism rather than a
heuristic, and no other statement in the paper does that. The bound that carries the
reported numbers is computed instead from the exact divergence between the two
release distributions at the operating tilt, and Section~2.9 places the two side by
side rather than letting the reader discover the gap.

\subsection{Privacy relative to a cohort mixture}

The baseline $R_D$ is the panel's own path model: the distribution over admissible
paths induced by the donor haplotypes embedded in the graph, with no reference to any
target. A draw from it is a mosaic of donor segments---a path the cohort could have
produced and that carries no information about the person being released. It is a
function of the graph and the panel alone, and the mechanism's output is a tilt of
it.

Privacy is stated against that object because it is the adversary's real
alternative. The pairwise statement of Section~2.6 answers ``which of these two
paths produced this release''. The question a cohort-scale release actually faces is
closer to ``is this anyone in particular'', and the natural alternative to the target
is not an arbitrary path but another member of the cohort, or the cohort taken as a
whole. That question has an answer of the same strength. Let $\pi$ be any prior over
the cohort and let
\begin{equation}
Q^{\pi,\tau}_{D} \;=\; \sum_{i=1}^{n} \pi_i\, Q^{\tau}_{p_i}
\end{equation}
be the induced mixture of cohort releases. Then
\begin{equation}
\mathrm{TV}\!\big(Q^{\tau}_{p_g},\, Q^{\pi,\tau}_{D}\big) \;\le\; \tau .
\label{eq:cor}
\end{equation}
This follows from the pairwise bound by convexity: total variation is convex in each
argument, so the distance from a fixed distribution to a mixture is at most the
largest distance to a component, and every component is within $\tau$ by
\eqref{eq:thm}. The bound therefore holds for every prior at once---the adversary's
own included---and does not require the adversary to disclose which prior it holds.

Both statements need the output support to be a single set fixed before the target is
seen, which is the first of the conditions in Section~2.8. The reason is not
technical: a release distribution supported on a target-dependent set is
distinguishable from a cohort member's on the support alone, and no bound on the
distance between the parts that do overlap repairs that.

The two regimes of Section~2.2 stand differently here, and the difference has to be
stated rather than smoothed over. For an external target the mixture runs over the
full panel and nothing in it depends on the target, so \eqref{eq:cor} applies as
written. For a cohort member it applies to the mixture over the panel with both of
that target's own haplotypes removed, which is why the leave-one-out construction of
Section~2.2 is a requirement of the guarantee and not a refinement of it; and even
then the graph itself was built with that target among its inputs, so in that regime
the guarantee is conditional on the graph. We state that limit and do not claim it
away.

\subsection{Conditions required by the guarantee}

The guarantee is not a property of the formula alone. It holds under three
conditions, and what makes them worth enumerating is that a release can violate any
of them and still produce entirely plausible output.

First, the output support and the baseline are built from the graph and the panel
alone, never from the target. In our implementation the only position excluded is one
at which no donor haplotype has a called allele, which is a property of the panel;
the support is then $R_D$'s own support, so no output the panel cannot produce is
released and none it can produce is withheld. Narrowing the support after the target
is known---to the paths near it, say---buys utility and voids the guarantee, and does
so without any visible symptom.

Second, the utility is bounded with a known range. Here $u$ maps into $[0,1]$ and
$\Delta u = 1$ because the weights sum to one, and the calibration of Section~2.3
divides by $\Delta u$. An unnormalised utility does not weaken the guarantee, it
removes it: the tilt is then unbounded in the output, and \eqref{eq:thm} has no
finite $\varepsilon_\tau$ to state.

Third, the release is sampled and never maximised. Maximum-a-posteriori decoding
returns the single most probable state sequence, which is not a draw from
$Q^{\tau}_{p}$ and inherits none of the bounds above, while being indistinguishable
from a legitimate release by inspection.

The guarantee also has a scope, and it is narrower than the person. It covers the
released path and nothing else: not the raw target genome, not the graph, and not any
further statistic published about the same individual. Two releases of one target are
two uses of the parameter, not one.

\subsection{The bound ladder}

Everything this paper bounds is a single quantity---how often an equal-prior
adversary holding one release and two candidate inputs can say which input produced
it---and by \eqref{eq:succ} that quantity is an affine function of one distance. What
differs between the statements in the paper is not the quantity but how conservatively
it is bounded, and the relationships between them are mathematics rather than
protocol choices. There are four, loosest first.

\begin{enumerate}
\item \textbf{The theorem.} $\mathrm{TV} \le \tau$ from \eqref{eq:thm}, hence
accuracy at most $(1+\tau)/2$. Proven for \emph{any} two admissible input paths, by
way of $\Delta u$, the full range of the utility function---a worst case the proof
has to survive.
\item \textbf{Pinsker from the exact divergence.} For an actual pair at the actual
tilt, with $D_{\mathrm{KL}} = D_{\mathrm{KL}}(Q^{\tau}_{p} \,\|\, Q^{\tau}_{q})$
computed exactly,
\begin{equation}
\mathrm{TV} \;\le\; \sqrt{D_{\mathrm{KL}}/2},
\qquad\text{hence}\qquad
P^{*}_{\mathrm{succ}} \;\le\; \tfrac{1}{2}\Big(1 + \sqrt{D_{\mathrm{KL}}/2}\Big).
\label{eq:pinsker}
\end{equation}
This is the quantity we report as the accuracy ceiling.
\item \textbf{The Gaussian conversion of the same divergence.} The privacy-loss
variable is a sum of millions of tiny per-position contributions and is Gaussian to
central-limit order, which gives
$\mathrm{TV} = 2\Phi\big(\sqrt{D_{\mathrm{KL}}/2}\big) - 1$ and hence
$P^{*}_{\mathrm{succ}} \le \Phi\big(\sqrt{D_{\mathrm{KL}}/2}\big)$.
\item \textbf{The best attacker we implemented.} This rung bounds the quantity from
\emph{below} and never from above. An attacker that fails is evidence about the
attacker.
\end{enumerate}

Rungs 1 to 3 are derivable from the column label alone, because \eqref{eq:pinsker}
inverts exactly: a reported ceiling $a$ corresponds to
$D_{\mathrm{KL}} = 2(2a-1)^2$. Rung 4 is measured and not derived.

Pinsker's conservatism relative to the Gaussian conversion can be quoted as measured
rather than as its limit. Taking the excess over chance on rung 2 divided by the
excess on rung 3, the ratio runs $1.2554$ at the $0.55$ operating point, $1.2721$ at
$0.65$ and $1.3289$ at $0.80$. The often-quoted $\sqrt{\pi/2} = 1.2533$ is the
small-divergence limit and understates the cost at the top of our grid, so we quote
the measured values.

The two routes point in opposite directions along two different axes, and conflating
them is what makes this section necessary. The divergence route is far tighter---it
yields $0.80$ where the theorem yields a bound of essentially $1$---but narrower in
scope: the theorem covers every input pair, while a ceiling covers the pair it was
computed for. On assumptions the ordering reverses, since the theorem needs only the
range of the utility whereas the divergence needs the model; given the model, though,
the divergence is exact, computed from the hidden Markov model with no Monte-Carlo
error rather than estimated from draws. The theorem is therefore more general and the
ceiling is more informative, and neither is simply the more robust of the two.

Because a ceiling is computed for one pair, a column label computed against one decoy
bounds a second pair only if the decoy it was computed against is the more divergent
of the two. We calibrate against the constructed extremal decoy of Section~2.11 and
then measure whether the resulting label dominates, rather than arguing that it must.
On chromosome 21 in the external-target regime it does, exactly and in every cell:
across ten targets, six tilted operating points and three real-genome decoy
conditions, the constructed decoy carries the larger exact divergence in all $180$
cells, and no real pair reaches half of it---the largest real-to-constructed
divergence ratio anywhere in that grid is $0.478$. In the reported units, the most
divergent real pair at the $0.80$ operating point has an exact ceiling of $0.708$
against the $0.800$ its column is labelled with. The same dominance holds in all
$264$ cells of the cohort-member regime on chromosome 21 and in all $60$ genome-wide
cells, for the decoy conditions measured at each scope. That is empirical dominance
over the pairs we can construct and not a proof over all of them, which is precisely
why rung 1 stays in the paper despite reading as a vacuous bound at every operating
point we report: it is the only statement covering every pair, and nothing else in
the paper does that job.

Two conventions follow, and the paper holds to both. The equivalent $\tau$ is stated
once, with its value, at the point the grid is introduced: at the $0.80$ operating
point it is within $10^{-2873}$ of one. Stating it defuses the obvious suspicion,
where omitting it invites a reader to derive it and wonder what else was left out.
And the ceiling is never called ``the guarantee'' without qualification. It is the
operational bound; the theorem is the structural one. Saying so once, early, lets the
rest of the paper write ``ceiling'' freely.
